\section{Learning the kernel}
\label{sec:laplace}

In practice the kernel is unknown and has to be learned from data. We work on the Laplace chain, a controlled case, but unlike the Gaussian, the score is nonlinear, and we compare learning the kernel with training neural denoisers via denoising score matching.

\subsection{Expectation--maximisation through BP}
\label{sec:em}

\begin{figure}[t]
  \centering
  \begin{tikzpicture}[
    box/.style={draw,thick,align=center,font=\small,minimum height=8.5mm,inner sep=2pt,text width=1.55cm},
    data/.style={box,fill=black!8},
    arr/.style={-{Stealth[length=2mm]},thick},
    lab/.style={font=\footnotesize,align=center,inner sep=1.5pt},
    stage/.style={font=\small\itshape,anchor=west,inner sep=0pt}]
  \node[stage] at (-0.85,1.2) {learn};
  \node[data] (x) at (0,0) {noisy\\ sequences};
  \node[box] (xi) at (3.05,0) {counts\\ $\Xi(a,a')$};
  \node[box] (k) at (6.1,0) {kernel\\ $K_\theta(a'|a)$};
  \draw[arr] (x) -- node[lab,above]{E-step} node[lab,below]{BP} (xi);
  \draw[arr] (xi) -- node[lab,above]{M-step} node[lab,below]{$\max_\theta Q$} (k);
  \draw[arr] (k.north) |- node[lab,above,pos=0.75]{repeat} (1.52,1.0) -- (1.52,0.52);
  \node[stage] at (-0.85,-1.5) {use};
  \node[data] (y) at (0,-2.55) {noisy\\ sequence $\mathbf{x}$};
  \node[box] (m) at (3.05,-2.55) {$\mathbb{E}[\mathbf{a}\mid\mathbf{x}]$};
  \node[box] (s) at (6.1,-2.55) {score\\ $\mathbf{S}_t(\mathbf{x})$};
  \draw[arr] (y) -- node[lab,above]{BP} node[lab,below]{any $t$} (m);
  \draw[arr] (m) -- node[lab,above]{Tweedie} (s);
  \draw[thick,dashed,-{Stealth[length=2mm]}] (k.south) |- (1.52,-1.45) -- (1.52,-2.05);
\end{tikzpicture}

  \caption{Learning and using the kernel. Learn: BP with the current kernel gives the posterior of every pair of
  neighbouring sites, summed into the expected transition counts $\Xi$ (E-step); the kernel is refitted to these
  counts (M-step); repeat. Use: BP with the learned kernel gives $\mathbb E[\mathbf a\mid\mathbf x]$, and so the score,
  at any noise level. BP enters the learning only through $\Xi$, so nothing is differentiated through BP.}
  \label{fig:em-loop}
\end{figure}

Let $K_\theta$ be a family of kernels and $\mathbf{x}^{(1)},\ldots,\mathbf{x}^{(N)}$ noisy training sequences, each
observed at a noise level $t_d$. The clean sequences are latent variables, and the likelihood of the data is
\begin{equation}
    p_\theta(\mathbf{x}^{(d)})
    =
    \int
    p_\theta(\mathbf{a})
    p_{t_d}(\mathbf{x}^{(d)}\mid\mathbf{a})
    \,d\mathbf{a}.
\end{equation}
This integral has no closed form, but BP computes it on the chain, from the normalising constants of the forward
messages (\cref{app:bp}).

We maximise it with expectation--maximisation (EM) \cite{dempster1977maximum}, as in the Baum--Welch algorithm for
hidden Markov models \cite{baum1970,rabiner1989tutorial}: the clean sequence is the hidden chain and the diffusion noise
is the emission (\cref{fig:em-loop}). The E-step uses the current parameters $\theta_{\mathrm{old}}$ to infer the
posterior of each pair of neighbouring sites,
\begin{equation}
    \xi_k^{(d)}(a,a')
    =
    p_{\theta_{\mathrm{old}}}
    \left(
        a_k=a,\,
        a_{k+1}=a'
        \mid \mathbf{x}^{(d)}
    \right),
\end{equation}
which the forward and backward messages give directly:
\begin{equation}
    \xi_k^{(d)}(a,a')
    \propto
    f_k(a)\,
    K_{\theta_{\mathrm{old}}}(a'\mid a)\,
    \ell_{t_d}(x_{k+1}^{(d)}\mid a')\,
    b_{k+1}(a').
    \label{eq:pairwise-posterior}
\end{equation}
Summing over sequences and links gives the expected transition counts
\begin{equation}
    \Xi(a,a')
    =
    \sum_{d=1}^{N}
    \sum_{k=1}^{L-1}
    \xi_k^{(d)}(a,a'),
\end{equation}
the expected number of transitions from $a$ to $a'$ in the clean data. The M-step then maximises
\begin{equation}
    Q(\theta\mid\theta_{\mathrm{old}})
    =
    \iint
    \Xi(a,a')\,
    \log K_\theta(a'\mid a)
    \,da\,da'
    \label{eq:em-objective}
\end{equation}
over $\theta$. With clean sequences nothing is hidden: $\Xi$ is just the histogram of the observed transitions and the
M-step is ordinary maximum likelihood. With noisy sequences the counts are soft, and BP provides them.

By Fisher's identity \cite{louis1982,cappe2005}, the gradient of $Q$ at $\theta_{\mathrm{old}}$ is
the gradient of the log-likelihood,
\begin{equation}
    \left.
    \nabla_\theta Q(\theta\mid\theta_{\mathrm{old}})
    \right|_{\theta=\theta_{\mathrm{old}}}
    =
    \nabla_\theta
    \log p_\theta(\mathbf{x})
    \big|_{\theta=\theta_{\mathrm{old}}},
\end{equation}
so each M-step moves uphill and the likelihood never decreases \cite{dempster1977maximum,wu1983}. Two practical
consequences follow. BP enters the M-step only through $\Xi$, which stays fixed while $\theta$ moves, so there is no
need to differentiate through the messages: the gradient only goes through the kernel and its normalisation. And the
noise level appears only in the E-step, so the kernel is learned once and then gives the score at any noise level.

\subsection{Kernel families}
\label{sec:kernel-families}

We first use families with the shift structure $a_{k+1}=\alpha a_k+\varepsilon_k$, where the noise $\varepsilon_k$ does
not depend on $a_k$. One option is a mixture of $C=8$ Gaussians for the noise,
\begin{equation}
    K_\theta(a'\mid a)
    =
    \sum_{c=1}^{C}
    \pi_c\,
    \mathcal{N}
    \left(
        a'-\alpha a;\nu_c,s_c^2
    \right),
\end{equation}
with weights $\pi_c$ summing to one. It has 25 parameters and an almost closed-form M-step, and it learns $\alpha$
quickly (to $0.1\%$ from 512 sequences), but it can only approximate the peak of the Laplace law by stacking Gaussians.
A second option writes the logarithm of the noise density as a truncated Fourier series, again with 25 parameters; this
keeps the kernel positive and makes the M-step concave. Other bases behave the same way (\cref{app:families}).
% src: thesis exp_18: alpha within 0.1% from N=512 (2.4% at 128)

Both families assume the shift structure. We therefore also use a kernel that assumes nothing beyond the Markov
property, a neural energy kernel, which we simply call the neural kernel:
\begin{equation}
    K_\theta(a'\mid a)
    =
    \frac{
        \exp[-E_\theta(a,a')]
    }{
        Z_\theta(a)
    },
    \label{eq:neural-kernel}
\end{equation}
with $Z_\theta(a)=\int\exp[-E_\theta(a,u)]\,du$ computed on the grid. $E_\theta$ is a small network with two hidden
layers of 16 units, 337 parameters in all, and the normalisation makes $K_\theta$ a valid transition density for any
weights.

With noisy sequences the neural kernel is trained by the same EM: BP gives $\Xi$ in the E-step, and the M-step takes
ten Adam steps \cite{kingma2015adam} on \cref{eq:em-objective}, differentiating only through the kernel. With clean
sequences it can instead be fitted by maximising the likelihood of the observed transitions,
\begin{equation}
    \mathcal{L}(\theta)
    =
    \sum_{d=1}^{N}
    \sum_{k=1}^{L-1}
    \log K_\theta
    \left(
        a_{k+1}^{(d)}
        \mid a_k^{(d)}
    \right),
    \label{eq:neural-clean-ml}
\end{equation}
with no E-step at all. The same kernel can therefore be learned from noisy sequences by EM or from clean ones by
maximum likelihood.

\subsection{Against neural denoisers}
\label{sec:neural-denoisers}

\begin{figure}[t]
  \centering
  \begin{tikzpicture}[
    box/.style={draw,thick,align=center,font=\small,minimum height=8.5mm,inner sep=2pt,text width=1.75cm},
    data/.style={box,fill=black!8,text width=2.3cm},
    arr/.style={-{Stealth[length=2mm]},thick},
    lab/.style={font=\footnotesize,align=center,inner sep=2pt}]
  \node[data] (a) at (0,0) {clean sequences};
  \node[box] (k) at (-1.9,-1.75) {kernel\\ $K_\theta(a'|a)$};
  \node[box] (n) at (1.9,-1.75) {network\\ $\boldsymbol{\epsilon}_\phi(\mathbf{x},t)$};
  \node[box] (s1) at (-1.9,-3.5) {score\\ $\widehat{\mathbf{S}}_t(\mathbf{x})$};
  \node[box] (s2) at (1.9,-3.5) {score\\ $\widehat{\mathbf{S}}_t(\mathbf{x})$};
  \draw[arr] (a.west) -| node[lab,left,pos=0.75]{max.\ likelihood} (k.north);
  \draw[arr] (a.east) -| node[lab,right,pos=0.75]{noise, DSM} (n.north);
  \draw[arr] (k) -- node[lab,left]{BP, any $t$} (s1);
  \draw[arr] (n) -- node[lab,right]{$-\boldsymbol{\epsilon}_\phi/\sqrt{\Delta_t}$} (s2);
\end{tikzpicture}

  \caption{Two ways to learn from the same clean sequences. Left: fit the kernel by maximum likelihood
  (\cref{eq:neural-clean-ml}); BP then gives the score at any noise level. Right: train a network by denoising score
  matching on noisy copies of the sequences (\cref{eq:dsm-objective}); it gives the score at the noise levels it was
  trained on. Only the kernel is estimated on the left: the score then follows exactly from BP, up to the grid and to the
  error of $\widehat K_\theta$. On the right the network estimates the score itself.}
  \label{fig:routes}
\end{figure}

We compare the learned kernels with neural denoisers trained by denoising score matching (DSM)
\cite{hyvarinen2005estimation,vincent2011connection}. These do not model the transition law; they learn the score
directly from noisy sequences. For a clean sequence $\mathbf a$ the noisy version is
\begin{equation}
    \mathbf{x}
    =
    e^{-t}\mathbf{a}
    +
    \sqrt{\Delta_t}\boldsymbol{\xi},
    \qquad
    \boldsymbol{\xi}
    \sim
    \mathcal{N}(\mathbf{0},\mathbf{I}),
\end{equation}
and a network $\boldsymbol{\epsilon}_\phi(\mathbf{x},t)$ is trained to predict the added noise by minimising
\begin{equation}
    \mathcal{L}_{\mathrm{DSM}}(\phi)
    =
    \mathbb{E}_{\mathbf{a},t,\boldsymbol{\xi}}
    \left[
        \left\|
            \boldsymbol{\xi}
            -
            \boldsymbol{\epsilon}_\phi
            (\mathbf{x},t)
        \right\|^2
    \right],
    \label{eq:dsm-objective}
\end{equation}
which gives the score estimate $\widehat{\mathbf{S}}_t(\mathbf{x})=-\boldsymbol{\epsilon}_\phi(\mathbf{x},t)/\sqrt{\Delta_t}$.

We use a global MLP, which knows nothing about locality, and a window CNN, which only looks at nearby sites (26{,}016
and 5{,}505 parameters). Both are trained on the clean sequences with fresh noise at every step (\cref{fig:routes}),
with hyperparameters chosen on separate seeds (\cref{app:networks}, which also has four other architectures). The
comparison is between two approaches: learning a small transition law and letting BP do the inference, or learning the
whole denoiser from examples. In the first the Markov structure is built into the model; in the second the network has
to find it in the data.
% src: outputs/E03/winners.json, confirm.csv n_params (mlp 26,016; window 5,505)

\subsection{Results on the Laplace chain}
\label{sec:results}

\begin{figure*}[t]
  \centering
  \includegraphics{fig_results}
  \caption{Relative score error against the number of training sequences $N$. (a) Laplace chain, 16 seeds.
  (b) Allen--Cahn field at $\Theta=0.3$ (\cref{sec:allen-cahn}), 8 seeds. Filled markers: kernels learned by EM from
  noisy sequences and used inside BP, the mixture of 8 Gaussians and the neural kernel. Open markers, dashed: the neural
  kernel fitted by maximum likelihood on the clean sequences. Global MLP and window CNN: denoising score matching on the
  clean sequences. Dotted: the best linear denoiser. Medians over seeds; on the Laplace chain the EM runs stop at
  $N=2048$ (\cref{app:em}).}
  \label{fig:results}
  % src: tools/paper_figs/figs_v2.py::fig_results. (a) E02a capacity.csv (em_bp, C=8), E56 kernels (energy), E58
  % (energy_h16, mlp, window, lmmse; E58's gate G2 reproduces E03's networks at N <= 2048). (b) E33 via read_e60
  % (em_mixture, raw_risk), E55 kernels (em_energy, selection=risk), E58 (energy_h16, mlp, window, lmmse); seeds 0-7.
\end{figure*}

We use the Laplace chain with $\alpha=0.85$ and $L=32$, and twelve noise levels between $t=0.05$ and $t=3$. Each model
is trained on $N$ sequences, from 32 to 2{,}048 (up to 32{,}768 with clean data), with 16 seeds. EM sees one noisy copy
of each training sequence, at one of the twelve levels; the networks see the clean sequences and draw fresh noise at
every step. We measure the relative score error $\|\widehat{\mathbf S}-\mathbf
S\|/\|\mathbf S\|$ on 2{,}048 test sequences, pooled over sites and all twelve levels (\cref{app:measures}), where $\mathbf S$ is the exact score from BP with the
true kernel. Every model keeps the checkpoint with the lowest error on separate validation sequences.

\Cref{fig:results}a shows the result. The mixture learned by EM, with 25 parameters, has the lowest error at every $N$,
from $0.099$ with 32 sequences to $0.023$ with 2{,}048. The global MLP is 6 to 10 times worse and the window CNN 2 to 6
times worse, in every seed. The CNN barely improves with more data and stays near $0.12$, not far from the best linear
denoiser ($0.17$): this is the locality limit of \cref{sec:gaussian}. The Fourier kernel is close to the mixture, 1.1 to
1.35 times its error. The neural kernel needs more data than the 25-parameter families, but it reaches the mixture with
2{,}048 sequences ($0.024$ against $0.023$), and it beats both networks in every seed at every $N$.
% src: E02a em_bp seed means 0.0989 ... 0.0231; per-seed mean ratios (E03, paired) MLP 5.79-9.85, window 2.07-5.58,
% 16/16 every N; window 0.1181 at 2048; E58 lmmse 0.1672; E56 fourier/mixture 1.08-1.35; energy 0.0235 at 2048,
% wins vs MLP and window 16/16 at every N

With clean sequences the neural kernel is fitted by maximum likelihood (open markers). It is then about three times
more accurate than from noisy sequences, and from only 128 clean sequences it beats the window CNN trained on 32{,}768.
(\cref{app:networks}).
% src: E58 (read_e58): E56 energy / E58 energy_h16 per-seed ratio 2.77-3.50; energy_h16 0.0251 at N=128 vs window 0.1162
% at N=32768; E66 (read_e66) P1: BP-DSM kernel below the best of six networks 16/16 at every N

\subsection{Step efficiency}
\label{sec:steps}

\begin{figure}[t]
  \centering
  \includegraphics{fig_steps}
  \caption{Step efficiency on the Laplace chain with $N=2048$: test score error against the number of parameter
  updates. One update is an EM iteration of the mixture, one Adam step of the neural kernel (ten per EM iteration), or
  one Adam step of a network on 64 sequences (\cref{app:measures}). Median over 16 seeds.}
  \label{fig:steps}
  % src: tools/paper_figs/figs_v2.py::fig_steps on reports/core/e63_curves_wall.csv (E65 refits, frozen 4 Nov 2026)
\end{figure}

The kernels also need far fewer parameter updates (\cref{fig:steps}). With 2{,}048 sequences the mixture passes the
final error of the window CNN after 20 EM iterations and is close to its own best after about 160. The neural kernel
passes the CNN's final error after 400 Adam steps from noisy data, and after 200 from clean data. The CNN needs all
of its 32{,}000 steps to get there.
% src: fig_steps printout (medians, E65): em_bp below 0.1176 at 20, within 10% of its minimum at 160; energy x10 below
% 0.1176 at 400; clean energy x10 at 200; window final 0.1176 at 32,000
Fewer updates are not always less time, because each EM iteration runs BP over all the training sequences. From clean
data the neural kernel is also faster in wall-clock time than both networks. (\cref{app:compute}).
% src: e63_curves_wall.csv at N = 2048: clean kernel first vs window and MLP 16/16; noisy kernel first 2/16

\subsection{When the data are not first-order Markov}
\label{sec:nonmarkov}

BP with a first-order kernel assumes first-order Markov data. When the data are not, the advantage over the networks
shrinks with the distance between the data and the closest first-order chain, and the networks win once this distance is relevant (\cref{app:nonmarkov}). For a second-order chain there is a simple fix: pairs of
neighbouring sites form a first-order chain again, so BP on pairs is exact, and EM on pairs brings the advantage back. This fix applies to higher-order chains, but at the cost of another dimension on the grid for each order.
% src: thesis exp_21_clean + E62: break-even delta 1-4e-2; E62 pair state below the best network >= 7/8 at every phi2
